\section{模型检验与实验设计}

\subsection{规则实现与工程验证}
为避免在正式测试前反复消耗机会，依据题面和附件中的公开规则搭建独立本地模拟器，对物理转移、接口协议、客户端重试、计时、回放和调试网页进行验证。当前可复现结果如\cref{tab:engineering-validation}所示。

\begin{table}[H]
  \centering
  \caption{当前工程验证结果}
  \label{tab:engineering-validation}
  \begin{tabularx}{\textwidth}{LYL}
    \toprule
    验证项目 & 结果 & 适用范围 \\
    \midrule
    自动化回归测试 & 59 项及 69 个参数化子测试通过 & 几何、覆盖、协议、回放与策略接口 \\
    附件计时回放 & $0,105,111,194,199,199$ s & 与附件给定动作序列一致 \\
    参考内核与快速内核对比 & 10000 步无差异 & 固定本地夹具 \\
    距离阈值邻域检查 & 48 组无差异 & 本地数值边界样例 \\
    题面及两份附件校验 & 3 份通过 & 仓库保存的原始文件 \\
    \bottomrule
  \end{tabularx}
\end{table}

上述测试说明已列出的公开规则和确定性夹具可复现，不代表本地场景与官方隐藏随机机制同分布，也不能代替问题三、问题四的完整策略效果实验。

\subsection{自建研究分布上的配对实验}
在固定的 \texttt{LOCAL-RESEARCH/research-v1-20260911} 分布上，对问题三的四组独立配置分别完成 512 次 PPO 更新，并分别在 64 个未参与更新的场景上与确定性教师策略配对比较，共得到 $256$ 对留出结果。两种策略的完成率均为 $100\%$；确定性基线与学习策略的平均虚拟时间分别为 $3743.90\,\mathrm{s}$ 和 $3353.06\,\mathrm{s}$，平均配对差为 $-390.85\,\mathrm{s}$，相对下降 $10.44\%$。按独立场景近似计算，平均差的 $95\%$ 区间半宽为 $45.29\,\mathrm{s}$。学习策略的平均路程由 $15131.36\,\mathrm{m}$ 降至 $13136.99\,\mathrm{m}$，平均失败清除次数由 $7.10$ 降至 $3.41$，但平均检测次数由 $105.97$ 略增至 $109.30$。

问题四此前仅有 $8$ 个小样本配对场景；虽然平均差为 $-114.60\,\mathrm{s}$，其近似 $95\%$ 区间跨越零。最新大规模训练记录进一步显示，四组教师基线的完成率仅为 $10.94\%$--$21.09\%$；重启后的 PPO 训练进行到第 154--157 次更新时，留出完成率仍为 $14.06\%$--$18.75\%$。因此问题四当前结果暴露的是终止与退出逻辑缺陷，不能作为策略改进证据，论文仍以连续域覆盖证书和确定性兜底作为可信结论。

\begin{table}[H]
  \centering
  \caption{自建研究分布上的策略配对结果}
  \label{tab:research-paired}
  \begin{tabularx}{\textwidth}{cYLLLL}
    \toprule
    问题 & 场景数 & 基线完成率 & 策略完成率 & 平均时间差/s & 近似 $95\%$ 半宽/s \\
    \midrule
    三 & $256$ & $100\%$ & $100\%$ & $-390.85$ & $45.29$ \\
    四（初步） & $8$ & $100\%$ & $100\%$ & $-114.60$ & $315.97$ \\
    \bottomrule
  \end{tabularx}
\end{table}

\begin{figure}[H]
  \centering
  \includegraphics[width=0.94\textwidth]{q3_large_paired.pdf}
  \caption{问题三 256 个留出场景的配对时间比较；虚线表示两策略耗时相同，负差值表示学习策略更快}
  \label{fig:q3-large-paired}
\end{figure}

训练与评估均只使用公开历史构造的状态特征；非法动作在采样前屏蔽。问题三最终配对验证的原始记录、配置和来源信息已从远端协作分支提交 \texttt{4d75bee} 原样归档至 \texttt{results/validation/q3-large-4d75bee/}，汇总值由 256 条记录重新计算。另以其中一组最优权重进行 3000 个独立随机场景的稳定性压力测试，全部完成任务，共清除 39024 个源；逐场景计算的平均每源耗时为 $261.65\,\mathrm{s}$，按全部源加权为 $256.23\,\mathrm{s}$。该试验没有配对基线，只支持稳定性判断；原始记录归档于 \texttt{results/validation/q3-3000-9e10994/}。以上均为非官方研究分布结果，正式运行仍以完成率不低于确定性基线为首要选择条件。

\subsection{策略评价指标}
每局记录实际源数 $N$、成功清除数 $C$、清除比例 $C/N$、累计虚拟时间 $T$、平均定位清除时间 $T/C$、程序运行时间、总路程、检测次数、频道切换次数、失败清除次数和最大单步决策耗时。比较方案时对相同场景进行配对测试，同时报告均值、标准差和失败案例。

\subsection{对照与灵敏度分析}
问题三至少比较 $a=1130,1135,1140,1150\,\mathrm{m}$ 的六边形覆盖骨架，以及“发现后立即追踪”和“覆盖途中顺路定位”两种调度。问题四比较不同网格平移、旋转和边长；任何候选首先通过连续域覆盖检查，再比较站点数和路线长度。问题二改变 $\lambda_1$、粗候选数量和局部细化范围，考察最坏后验半径与总时间的权衡。\placeholder{补充问题二参数灵敏度图，以及问题三、四的完整轨迹与退出证书图；新增数据必须链接到固定代码版本及原始运行记录。}
